%------------------------------------------------------------------------------%
%------------------------------------------------------------------------------%
\section{Integral Indefinida}

Na subseção anterior fizemos o processo inverso do que vinhamos fazendo até
aqui. Antes, tinhamos uma função $f$ e gostaríamos de encontrar a sua derivada,
agora dada uma determinada função $f'$ estamos interessado em encontrar uma
função (na verdade, um família de funções) $f$ cuja a derivada seja $f'$.

Em síntese temos que $f'$ é derivada de $f \Leftrightarrow f$ é uma primitiva
de $f'$.

A integral indefinida de uma função pode ser vista exatamente como a família de
primitivas da função.

%------------------------------------------------------------------------------%
\begin{definicao}{Integral Indefinida}{def:integral-indefinida}
  A notação $\int f(x) dx$ é tradicionalmente usada para a
  primitiva de $f$ e é chamada de integral indefinida, ou seja
  $\int f(x) dx= F(x)$ significa que $F$ é a primitiva mais geral de $f$,
  ou seja, $F'(x)=f(x)$.
\end{definicao}
%------------------------------------------------------------------------------%

Decorre das propriedades de primitiva que a integral indefinida preserva as
propriedades de linearidade, isto é,
\begin{itemize}
  \item \(\int f(x)+g(x) dx = \int f(x) dx + \int g(x) dx\),
  \item \(\int c f(x) dx = c \int f(x) dx\).
\end{itemize}

%------------------------------------------------------------------------------%
\begin{example}~
  \solution
  \begin{itemize}
    \item \(\int x^2 dx = \frac{x^3}{3}+C\)
    pois \(\frac{d}{dx} \left(  \frac{x^3}{3}+C \right)=x^2\)
    \item \(\int sec^2(x) dx = \tan(x)+C \)
    pois \(\frac{d}{dx} \left( \tan(x)+C \right)=\sec^2(x)\).
  \end{itemize}
\end{example}
%------------------------------------------------------------------------------%

%------------------------------------------------------------------------------%
\begin{example}
  Encontre a integral indefinida geral $\int (10x^4-2sec^2(x))dx$.

  \solution
  Usando as propriedades de integral indefinida
  \begin{align*}
    \int (10x^4-2sec^2(x))dx
    & = 10 \int x^4 dx - 2 \int \sec^2(x)dx \\
    & = 10 \frac{x^5}{5}-2\tan(x)+C \\
    & = 2x^5-2\tan(x)+C
  \end{align*}
\end{example}
%------------------------------------------------------------------------------%

%------------------------------------------------------------------------------%
\begin{example}
  Uma partícula move-se em uma reta e tem aceleração dada por $a(t)=6t+4$. Sua
  velocidade inicial é $v(0)=6 cm/s$, e seu deslocamento inicial é $s(0)=9cm$.
  Encontre sua função posição.

  \solution
  Uma vez que $v'(t)=a(t)=6t+4$, então:
  \[
  v(t)= \int a(t) dt= 3t^2+4t+C
  \]
  Podemos encontrar o valor da constante $C$ pois nos foi dado que $v(0)=6$.
  Substituindo essa informação na equação da velocidade encontrada
  \begin{align*}
    6 & = 3,0^2 + 4.0 + C \\
    C & =-6
  \end{align*}
  Assim, $v(t)=3t^2+4t-6$. Sabendo que $s'(t)=v(t)$, então $s$
  é a primitiva de $v$
  \[
  s(t)=\int v(t) dt
  = \int 3t^2+4t-6 dt
  = t^3+2t^2-6t+D
  \]
  Sabendo que $s(0)=9$, , temos que $D=9$. Portanto, a função posição pedida é
  \[
  s(t) = t^3+2t^2-6t+9.
  \]
\end{example}
%------------------------------------------------------------------------------%

